% GENERATED FILE. Edit canonical lessons, then run npm run generate:books.
% canonical-source-hash: fnv1a64:563e8fbb612975c9
% canonical-lessons: ZH-C149-zhengzai, ZH-W149-zheng, ZH-C149-zai-progressive, ZH-C149-zhenghao, ZH-C149-yibian, ZH-C149-practice

\chapter{Right Now}
\label{ch:zh-right-now}
\hlchaptermodality{149}

\begin{chapteropening}
\textbf{By the end of this chapter:} \emph{I can say what someone is doing right now.}

Saying what is going on right now with \zh{在} and \zh{正在} before a verb, two things at once with \zh{一边} … \zh{一边}, and a happy coincidence with \zh{正好}, built on one new written character: \zh{正}.
\end{chapteropening}

\section[\texorpdfstring{\zh{正在} (zhèngzài)}{正在 (zhèngzài)}]{\zh{正在} (zhèngzài) — in the middle of (doing)}

\label{lesson:ZH-C149-zhengzai}



\begin{quote}
A new chapter, and it opens on what is going on at this moment.

\begin{itemize}
\raggedright
  \item \emph{Recall:} say \zh{下个月我要去中国} — \emph{xià ge yuè wǒ yào qù Zhōngguó} — \textbf{R1}, the chapter you have finished
  \item \emph{Recall:} say \emph{jiānglái}, in the future — \textbf{R1}
  \item \emph{Recall:} say \zh{我在家} — \emph{wǒ zài jiā}, I'm at home — \textbf{R3}, durable
\end{itemize}
\end{quote}

\subsection*{You'll want to know: \zh{正在}}
\textbf{\zh{正在}} — \emph{zhèngzài} — \textquotedblleft{}in the middle of (doing)\textquotedblright{}.

\begin{quote}
\textbf{\zh{我正在吃饭}} \emph{wǒ zhèngzài chī fàn} — I'm eating right now.
\end{quote}

\begin{quote}
\textbf{\zh{他正在看书}} \emph{tā zhèngzài kàn shū} — He's reading at the moment.
\end{quote}

The next lesson puts the first character in your hand.

\begin{cousinweb}
\textbf{\zh{正在}} (\emph{zhèngzài}) is built from \zh{正} 'exactly' and \zh{在} 'at', the \zh{在} you write. Both are inherited from Middle Chinese, \emph{tsyengH} and \emph{dzojX} (Guangyun, in Baxter's transcription), glossed 'exactly fitting; proper' and 'to dwell; to exist'. The use before a verb, in the middle of doing it, is the modern one and is assumed here, not quoted.
\end{cousinweb}

\begin{grammarlens}[title={What you've built}]
\zh{正在} goes before the verb: \zh{我}, \zh{正在}, \zh{吃饭}.

\zh{在} without \zh{正} can do the same job before a verb; that comes two lessons on.
\end{grammarlens}

\subsection*{Your turn}
\begin{itemize}
\raggedright
  \item \emph{Say it:} \emph{zhèngzài}
  \item \emph{Say it:} \emph{zhèngzài}, once more
  \item \emph{Say it:} \zh{我正在吃饭} — \emph{wǒ zhèngzài chī fàn} — I'm eating right now.
  \item \emph{Say it:} \zh{他正在看书} — \emph{tā zhèngzài kàn shū} — He's reading at the moment.
  \item \emph{Recall:} say \zh{我在家}, then \zh{我正在吃饭}
\end{itemize}

\begin{tcolorbox}[breakable,colback=teal!4,colframe=teal!35!black,title={Before you move on}]
What does \zh{正在} mean? (\textquotedblleft{}In the middle of (doing)\textquotedblright{}.) Where does it go? (\textbf{Before the verb}.)
\end{tcolorbox}
\section[\texorpdfstring{\zh{正} (zhèng)}{正 (zhèng)}]{\zh{正} — \zh{一} on top, then the middle of \zh{上} on a long bar}

\label{lesson:ZH-W149-zheng}



\begin{quote}
The word first, then the character inside it.

\begin{itemize}
\raggedright
  \item \emph{Recall:} say \emph{zhèngzài} — \textbf{R1}, one lesson back, before it decays
\end{itemize}

\begin{quote}
\zh{一} \zh{上}
\end{quote}

\emph{Read it:} each one aloud, with nothing said first — \emph{yī}, \emph{shàng}
\end{quote}

\subsection*{Script — \zh{一} on top, then the middle of \zh{上} on a long bar}
\hlblockfigure{\detokenize{figures/ZH-W149-zheng-filmstrip.pdf}}{How \zh{正} is written, stroke by stroke}

\begin{quote}
\zh{正}
\end{quote}

\textbf{\zh{一} on top, then four strokes, two of them shaped as in \zh{上}.}

\textbf{\zh{一}}, a bar across the top.

\textbf{A vertical down the middle}, from under the bar to the foot.

\textbf{A short bar to its right}, halfway down: with the vertical, the upper part of \zh{上}.

\textbf{A short vertical on the left}, from the height of that bar down to the foot.

\textbf{The long bar along the foot}, last, under everything.

1 + 4 = five.

\subsection*{Your turn}
\begin{itemize}
\raggedright
  \item \emph{Write it:} \zh{正}, all five, with the model in front of you
  \item \emph{Cover:} write it again from memory, and count the pieces, not the strokes
\end{itemize}

\begin{tcolorbox}[breakable,colback=teal!4,colframe=teal!35!black,title={Before you move on}]
How many strokes in \zh{正}? (\textbf{Five}.) Which stroke is last? (\textbf{The long bar along the foot}.) Which character you write shares its middle? (\textbf{\zh{上}}.)

Source: \href{https://raw.githubusercontent.com/chanind/hanzi-writer-data/68d10a4b21150cae5e1ebbd223eed289cf32d90c/data/\%E6\%AD\%A3.json}{Hanzi Writer Data \zh{正}, pinned stroke paths and medians}.
\end{tcolorbox}
\section[\texorpdfstring{\zh{你在做什么} (nǐ zài zuò shénme)}{你在做什么 (nǐ zài zuò shénme)}]{\zh{你在做什么} — \zh{在} before a verb}

\label{lesson:ZH-C149-zai-progressive}



\begin{quote}
\zh{在} again, and this time it comes before a verb.

\begin{itemize}
\raggedright
  \item \emph{Read it:} \zh{正在}, aloud — \emph{zhèngzài}
  \item \emph{Recall:} say \zh{你呢} — \emph{nǐ ne}, and you?
\end{itemize}
\end{quote}

\begin{grammarlens}[title={Grammar Lens: \zh{在} before a place, \zh{在} before a verb}]
\begin{quote}
\textbf{\zh{在} before a place says where. \zh{在} before a verb says what is going on right now.}
\end{quote}

\begin{quote}
\textbf{\zh{他在家}} \emph{tā zài jiā} — He's at home. \textbf{\zh{他在看书}} \emph{tā zài kàn shū} — He's reading.
\end{quote}

\zh{正在} says it more firmly, in the very middle of it: \zh{我正在吃饭}. \zh{呢} at the end can round off either one:

\begin{quote}
\textbf{\zh{他在看书呢}} \emph{tā zài kàn shū ne} — He's busy reading.
\end{quote}

The negative is \zh{没在}, since \zh{没} reports what is not going on:

\begin{quote}
\textbf{\zh{我没在看书}} \emph{wǒ méi zài kàn shū} — I'm not reading.
\end{quote}

\begin{quote}
\textbf{\zh{你在做什么}} \emph{nǐ zài zuò shénme} — What are you doing?
\end{quote}
\end{grammarlens}

\subsection*{Your turn}
Say these aloud:

\begin{itemize}
\raggedright
  \item \zh{你在做什么} — \emph{nǐ zài zuò shénme} — What are you doing?
  \item \zh{我在看书} — \emph{wǒ zài kàn shū} — I'm reading.
  \item \zh{他在看书呢} — \emph{tā zài kàn shū ne} — He's busy reading.
  \item \zh{我没在看书} — \emph{wǒ méi zài kàn shū} — I'm not reading.
\end{itemize}

\begin{tcolorbox}[breakable,colback=teal!4,colframe=teal!35!black,title={Before you move on}]
What does \zh{在} say before a place? (\textbf{Where}.) Before a verb? (\textbf{What is going on now}.) How do you say you are not reading? (\textbf{\zh{我没在看书}}.)
\end{tcolorbox}
\section[\texorpdfstring{\zh{正好} (zhènghǎo)}{正好 (zhènghǎo)}]{\zh{正好} (zhènghǎo) — exactly right; as it happens}

\label{lesson:ZH-C149-zhenghao}



\begin{quote}
Before the new one: read the last lesson's question, then the word for good.

\begin{itemize}
\raggedright
  \item \emph{Read it:} \zh{你在做什么}, aloud — \emph{nǐ zài zuò shénme}
  \item \emph{Recall:} say \zh{好} — \emph{hǎo}, good — \textbf{R4}, recognition at distance
\end{itemize}
\end{quote}

\subsection*{You'll want to know: \zh{正好}}
\textbf{\zh{正好}} — \emph{zhènghǎo} — \textquotedblleft{}exactly right; as it happens\textquotedblright{}.

\begin{quote}
\textbf{\zh{正好} \zh{我也要去}} \emph{zhènghǎo, wǒ yě yào qù} — As it happens, I'm going too.
\end{quote}

\begin{quote}
\textbf{\zh{这个正好}} \emph{zhège zhènghǎo} — This one is exactly right.
\end{quote}

\begin{cousinweb}
\textbf{\zh{正好}} (\emph{zhènghǎo}) is built from \zh{正} 'exactly' and \zh{好} 'good': exactly right. \zh{好} is inherited from Middle Chinese \emph{xawX} (Guangyun, in Baxter's transcription), glossed 'good; fine'.
\end{cousinweb}

\begin{grammarlens}[title={What you've built}]
\zh{正好} can open a sentence, a happy coincidence: \zh{正好}, \zh{我也要去}.

Or it can follow a thing, which then fits: \zh{这个正好}.
\end{grammarlens}

\subsection*{Your turn}
\begin{itemize}
\raggedright
  \item \emph{Say it:} \emph{zhènghǎo}
  \item \emph{Say it:} \emph{zhènghǎo}, once more
  \item \emph{Say it:} \zh{正好} \zh{我也要去} — \emph{zhènghǎo, wǒ yě yào qù} — As it happens, I'm going too.
  \item \emph{Say it:} \zh{这个正好} — \emph{zhège zhènghǎo} — This one is exactly right.
  \item \emph{Recall:} say \zh{你在做什么}, then \zh{正好}
\end{itemize}

\begin{tcolorbox}[breakable,colback=teal!4,colframe=teal!35!black,title={Before you move on}]
What does \zh{正好} mean? (\textquotedblleft{}Exactly right; as it happens\textquotedblright{}.) Say that, as it happens, you are going too. (\textbf{\zh{正好} \zh{我也要去}}.)
\end{tcolorbox}
\section[\texorpdfstring{\zh{一边} (yìbiān)}{一边 (yìbiān)}]{\zh{一边} (yìbiān) — at the same time}

\label{lesson:ZH-C149-yibian}



\begin{quote}
Before the new one: last lesson's word, then a word with the same ending.

\begin{itemize}
\raggedright
  \item \emph{Recall:} say \zh{正好} — \emph{zhènghǎo}, as it happens — \textbf{R1}, one lesson back
  \item \emph{Recall:} say \emph{pángbiān}, beside — \textbf{R3}, durable
\end{itemize}
\end{quote}

\subsection*{You'll want to know: \zh{一边}}
\textbf{\zh{一边}} — \emph{yìbiān} — \textquotedblleft{}at the same time\textquotedblright{}, said twice, before two verbs; also \textquotedblleft{}one side\textquotedblright{}.

\begin{quote}
\textbf{\zh{他一边吃饭一边看书}} \emph{tā yìbiān chī fàn yìbiān kàn shū} — He reads while he eats.
\end{quote}

\zh{一} falls before the level tone of biān: \emph{yìbiān}, as in \emph{yì tiān}.

\begin{cousinweb}
\textbf{\zh{一边}} (\emph{yìbiān}) is built from \zh{一} 'one' and \zh{边} 'side', the \zh{边} you write. Both are inherited from Middle Chinese, \emph{'jit} and \emph{pen} (Guangyun, in Baxter's transcription), \emph{pen} glossed 'an edge; a border'.
\end{cousinweb}

\begin{grammarlens}[title={What you've built}]
Said twice, each straight before its own verb, \zh{一边} … \zh{一边} … puts two actions side by side: \zh{一边吃饭}, \zh{一边看书}.

The person is said once, at the front: \zh{他一边吃饭一边看书}.
\end{grammarlens}

\subsection*{Your turn}
\begin{itemize}
\raggedright
  \item \emph{Say it:} \emph{yìbiān}
  \item \emph{Say it:} \emph{yìbiān}, once more
  \item \emph{Say it:} \zh{他一边吃饭一边看书} — \emph{tā yìbiān chī fàn yìbiān kàn shū} — He reads while he eats.
  \item \emph{Say it:} \emph{wǒ yìbiān hē chá yìbiān shuō huà} — I talk while I drink my tea.
  \item \emph{Recall:} say \emph{pángbiān}, then \emph{yìbiān}
\end{itemize}

\begin{tcolorbox}[breakable,colback=teal!4,colframe=teal!35!black,title={Before you move on}]
What does \zh{一边} … \zh{一边} do? (\textbf{Puts two actions at the same time}.) Where does each \zh{一边} go? (\textbf{Straight before its verb}.)
\end{tcolorbox}
\section[Practice]{\zh{我正在吃饭}}

\label{lesson:ZH-C149-practice}



\begin{quote}
The chapter closes here, with a phone call about what everyone is doing.

\begin{itemize}
\raggedright
  \item \emph{Recall:} say \zh{一边吃饭一边看书} — \emph{yìbiān chī fàn yìbiān kàn shū} — \textbf{R1}, one lesson back
  \item \emph{Recall:} say \zh{正好} — \emph{zhènghǎo} — \textbf{R1}
  \item \emph{Recall:} say \emph{bànyè}, the middle of the night — \textbf{R2}, the chapter on night
\end{itemize}
\end{quote}

\subsection*{The exchange}
\begin{quote}
\textbf{A:} \zh{你在做什么}
\end{quote}

\begin{quote}
\emph{nǐ zài zuò shénme} — What are you doing?
\end{quote}

\begin{quote}
\textbf{B:} \zh{我正在吃饭} \zh{你呢}
\end{quote}

\begin{quote}
\emph{wǒ zhèngzài chī fàn, nǐ ne} — I'm eating. And you?
\end{quote}

\begin{quote}
\textbf{A:} \zh{我在看书}
\end{quote}

\begin{quote}
\emph{wǒ zài kàn shū} — I'm reading.
\end{quote}

\begin{quote}
\textbf{B:} \zh{你妈妈呢}
\end{quote}

\begin{quote}
\emph{nǐ māma ne} — And your mum?
\end{quote}

\begin{quote}
\textbf{A:} \zh{妈妈在做饭}
\end{quote}

\begin{quote}
\emph{māma zài zuò fàn} — She's cooking.
\end{quote}

\begin{quote}
\textbf{B:} \emph{nǐ bàba ne} — And your dad?
\end{quote}

\begin{quote}
\textbf{A:} \zh{他没在家} \zh{他在开车}
\end{quote}

\begin{quote}
\emph{tā méi zài jiā, tā zài kāichē} — He isn't home. He's driving.
\end{quote}

\begin{quote}
\textbf{B:} \zh{我一边吃饭一边打电话}
\end{quote}

\begin{quote}
\emph{wǒ yìbiān chī fàn yìbiān dǎ diànhuà} — I'm eating and on the phone at once.
\end{quote}

\begin{quote}
\textbf{A:} \zh{正好} \zh{我也在吃饭}
\end{quote}

\begin{quote}
\emph{zhènghǎo, wǒ yě zài chī fàn} — As it happens, I'm eating too.
\end{quote}

Every answer keeps \zh{在} before its verb. \zh{他没在家} is \zh{在} before a place, at home; \zh{他在开车} is \zh{在} before a verb, driving.

\begin{tcolorbox}[breakable,skin=enhanced,colback=orange!6,colframe=orange!45!black,boxrule=0.5pt,arc=1mm,left=6pt,right=6pt,top=4pt,bottom=4pt,fonttitle=\bfseries,title={Read this}]
\emph{Read it:} the exchange through once without stopping, then answer the question below
\end{tcolorbox}

\begin{grammarlens}[title={What you've built — now, and two things at once}]
\zh{在} or \zh{正在} before the verb says what is going on now:

\begin{itemize}
\raggedright
  \item \textbf{\zh{我正在吃饭}} \emph{wǒ zhèngzài chī fàn} — I'm eating right now.
  \item \textbf{\zh{他在开车}} \emph{tā zài kāichē} — he's driving.
  \item \textbf{\zh{白天我在看书}} \emph{báitiān wǒ zài kàn shū} — in the daytime I'm reading.
  \item \emph{bànyè tā zhèngzài shuìjiào} — in the middle of the night he's asleep.
\end{itemize}

\textbf{Listen.} \emph{Nǐ zài zuò shénme ne?}

\textbf{Speak.} Say what you are doing now, then what somebody at home is doing.

\emph{Write it:} \zh{正在} — \zh{一}, then the middle of \zh{上}, then the long bar
\end{grammarlens}

\subsection*{Your turn}
\begin{itemize}
\raggedright
  \item \emph{Say it:} \zh{你在做什么} — \emph{nǐ zài zuò shénme} — What are you doing?
  \item \emph{Say it:} \zh{我正在吃饭} \zh{你呢} — \emph{wǒ zhèngzài chī fàn, nǐ ne} — I'm eating. And you?
  \item \emph{Say it:} the whole exchange, both parts
  \item \emph{Write it:} \zh{我正在吃饭}
\end{itemize}

\begin{tcolorbox}[breakable,colback=teal!4,colframe=teal!35!black,title={Before you move on}]
Ask what someone is doing. (\textbf{\zh{你在做什么}}.) Say you are eating right now. (\textbf{\zh{我正在吃饭}}.) Say he reads while he eats. (\textbf{\zh{他一边吃饭一边看书}}.)
\end{tcolorbox}
